\section{The Simulation-Based Inference Paradigm}
\label{sec:3.2}

The inference framework presented in Section~\ref{sec:3.1} assumes that the likelihood function can be evaluated explicitly.
This assumption holds for many standard analyses, but it breaks down when the data-generating process involves a complex simulator whose internal stochastic structure prevents the likelihood from being written in closed form.
Simulation-based inference (SBI) provides a principled alternative: \blue{when evaluating the likelihood analytically is impossible or impractical}, one trains machine learning models on simulated data to approximate the posterior, the likelihood, or their ratio.
This section introduces the SBI paradigm, the neural estimation strategies \blue{behind} it, and the validation techniques required to ensure the reliability of the results.

%% --- 3.2.1 ---

\subsection{The Intractable Likelihood Problem}
\label{sec:3.2.1}

The inference methods presented in Section~\ref{sec:3.1} share a common requirement: the ability to evaluate the likelihood function $p(\mathbf{x}|\boldsymbol{\theta})$ for a given set of parameters.
In many problems of interest, however, the data-generating process is described by a complex computer simulation, \blue{and the results of such simulations are formally samples from the implicit likelihood.}
Such models are known as \emph{implicit models}, in contrast to \emph{prescribed models} for which the likelihood can be written in closed form~\cite{Cranmer:2020,Diggle:1984}.
The resulting inference problem has historically been called \emph{likelihood-free inference}, although the term is somewhat misleading: these methods typically estimate the intractable likelihood rather than dispensing with it entirely.
\blue{The more descriptive term \emph{implicit likelihood methods} is now preferred~\cite{Cranmer:2020}, and \emph{simulation-based inference} (SBI) is a term often used as an equally valid synonim. For convenience, and to align with the machine learning literature, we will refer to simulation-based inference throughout this thesis.}

The source of the difficulty is straightforward to state.
A typical simulator takes as input a vector of physical parameters $\boldsymbol{\theta}$, samples a series of internal stochastic states or latent variables $z_i$,\footnote{\blue{In the context of implicit likelihood methods, the latent variables $z_i$ represent the internal stochastic states of the simulator that are not directly observable. For example, in the context of the work of Chapter \ref{ch:6}, that would be the position of the indibidual gamma-ray sources.}} and produces synthetic data $\mathbf{x}$ as output.
The exact likelihood of this output requires marginalizing over the joint density of data and latent variables:
\begin{equation}
	\label{eq:intractable_likelihood}
	p(\mathbf{x}|\boldsymbol{\theta}) = \int p(\mathbf{x}, \mathbf{z}|\boldsymbol{\theta}) \, d\mathbf{z},
\end{equation}
where the integral runs over every possible latent variable configuration of the simulator~\cite{Cranmer:2020}.
For realistic simulations involving billions of latent variables, such as particle transport through a detector, the stochastic evolution of astrophysical populations, or the formation of large-scale structure, this high-dimensional integral is computationally intractable.

The asymmetry between the forward and inverse problems is a defining feature of the SBI setting.
Running the simulator to produce synthetic observations from a given $\boldsymbol{\theta}$ is computationally doable: one simply executes the program.
Inverting this process to obtain $p(\mathbf{x}|\boldsymbol{\theta})$, however, would require enumerating all possible internal configurations that could have led to the same observable output.
This forward--inverse asymmetry motivates an alternative strategy: rather than evaluating the likelihood directly, one can generate large numbers of simulated datasets and use them to learn the relationship between parameters and observables.

A concrete astrophysical example arises in the measurement of the gamma-ray source-count distribution $dN/dS$ below the Fermi-LAT detection threshold.
As discussed in detail in Chapter~\ref{ch:6}, the observed photon-count map results from the superposition of thousands of unresolved point sources.
Writing a closed-form likelihood for this full sky map would require integrating over every possible spatial configuration, flux, and number of the unresolved sources.
Generating a synthetic map from a given set of $dN/dS$ parameters, by contrast, amounts to a sequence of well-defined sampling steps: draw source fluxes from the $dN/dS$, scatter them on the sky, and apply additional well-understood sources of noise.
The forward simulation is straightforward; the inverse problem is not.
SBI addresses exactly this kind of problem, and the following subsections describe how the inversion is accomplished.

%% --- 3.2.2 ---

\subsection{Simulation-Based Inference}
\label{sec:3.2.2}

Given the impossibility of evaluating the likelihood directly, SBI replaces explicit likelihood computation with learned surrogates trained on simulated data.
The idea has a long history: Diggle and Gratton~\cite{Diggle:1984} proposed using kernel density estimation to approximate intractable likelihoods from simulator output as early as 1984, while Approximate Bayesian Computation (ABC) developed in parallel as an entirely sampling-based alternative (see~\cite{Cranmer:2020} for a review).
Both approaches, however, suffer from the curse of dimensionality: they require hand-crafted summary statistics to compress the data into low-dimensional representations, inevitably discarding information about the parameters of interest.
Modern neural SBI methods overcome these limitations by using deep neural networks as flexible density estimators that potentially operate directly on high-dimensional data~\cite{Cranmer:2020}.
Three main neural strategies have emerged, distinguished by which statistical quantity the neural network learns to approximate.

\emph{Neural Posterior Estimation} (NPE) trains a conditional density estimator---\blue{typically a discrete or continuous normalizing flow~\cite{Papamakarios:2021,Lipman:2023hef}}---to directly approximate the posterior distribution $p(\boldsymbol{\theta}|\mathbf{x})$~\cite{Greenberg:2019}.
By conditioning the density estimator on the observed data $\mathbf{x}$, the network learns to output the posterior density for any input observation.
Because NPE directly provides the posterior, it is possibly the most straightforward approach for Bayesian inference.
However, because the prior enters the training, changing the prior requires retraining.

\emph{Neural Likelihood Estimation} (NLE) instead approximates the likelihood function $p(\mathbf{x}|\boldsymbol{\theta})$, conditioning the density estimator on $\boldsymbol{\theta}$ rather than on $\mathbf{x}$~\cite{Cranmer:2020}.
Once the likelihood surrogate is trained, \blue{Bayesian inference proceeds through standard means, for example through MCMC sampling with any choice of prior.}
This separation of the likelihood from the prior provides additional flexibility, at the cost of an extra sampling step.

\emph{Neural Ratio Estimation} (NRE) takes a different approach altogether: rather than learning a probability density, a classifier is trained via supervised learning to discriminate between synthetic datasets generated at different parameter points. \blue{This is attractive because classification is sometimes a more tractable learning problem than density estimation, especially in high-dimensional spaces.}
The classifier output can be converted into the likelihood-to-evidence ratio $r(\mathbf{x},\boldsymbol{\theta}) = p(\mathbf{x}|\boldsymbol{\theta})/p(\mathbf{x})$, which is proportional to the posterior up to the prior~\cite{Cranmer:2020,Miller2021arxiv}.
% A key advantage of NRE is its ability to target arbitrary low-dimensional marginal posteriors directly, without first estimating the full joint posterior~\cite{Cranmer:2020,Miller2021arxiv}.
% This makes the approach particularly attractive when one- and two-dimensional marginals suffice, as is common in cosmological and astrophysical parameter estimation.
% A key advantage of NRE is its ability to target arbitrary low-dimensional marginal posteriors directly, without first estimating the full joint posterior on both the parameters of interest and the nuisance parameters~\cite{Cranmer:2020,Miller2021arxiv}.
Once the classifier is trained, evaluating the posterior density at any point is fast, but drawing samples still requires MCMC. \blue{This is possibly the main drawback of NRE, as full posterior sampling and calibration are significantly more expensive than in NPE.}

\blue{All three strategies provide \emph{amortized inference}: the computationally expensive simulations are performed once. 
Training a model does not require running the simulator for every new observation, and different model architectures can be explored without additional calls to the simulator.
Once a model has been trained, it can be applied to sample the posterior for any new observed data, without additional calls to the simulator~\cite{Cranmer:2020}.
The neural network also implicitly marginalizes over the vast latent space of the simulator, so the trained surrogate is likely to preserve information that hand-crafted summary statistics would discard --- which makes SBI well suited to field-level inference. }

\paragraph{Application preview.}

In this thesis, NPE is the SBI strategy most directly relevant to our work.
In Chapter~\ref{ch:6}, a convolutional neural network trained on synthetic gamma-ray sky maps estimates the source-count distribution $dN/dS$ and its uncertainty for any input map, without evaluating an explicit likelihood.
The approach operates in the spirit of NPE: the network directly maps an observed sky map to posterior summaries (the mean and variance for each flux bin).
% , though it uses a heteroscedastic Gaussian loss (Section~\ref{sec:3.1.3}) rather than a normalizing flow to parameterize the output distribution.
\blue{Once trained, the inference is amortized: the network can be applied to the real Fermi-LAT sky map to quickly obtain the source-count distribution without requiring new simulations.}

%% --- 3.2.3 ---

\subsection{Validation of Simulation-Based Inference}
\label{sec:3.2.3}

A neural density estimator trained on finite simulations can, in principle, produce posteriors that are poorly calibrated: the reported credible intervals may be systematically too narrow or too wide.
Validating SBI results is therefore an integral part of any simulation-based analysis, and the simulator itself provides the means for doing so. 
\blue{We notice that validation would in principle also be required even when performing inference with a known likelihood, but it is often overlooked in practice as it is often computationally intractable for standard MCMC methods.}

\blue{The standard validation strategy is \emph{coverage testing}, and the procedure is conceptually simple.}
One samples parameters $\boldsymbol{\theta}^*$ from the prior, generates synthetic data $\mathbf{x}^* \sim p(\mathbf{x}|\boldsymbol{\theta}^*)$ using the simulator, and then applies the trained estimator to obtain the approximate posterior $\hat{p}(\boldsymbol{\theta}|\mathbf{x}^*)$.
Because the true parameters are known by construction, one can check whether the posterior places the correct amount of probability mass around $\boldsymbol{\theta}^*$.
Repeating this test across hundreds or thousands of independent simulations yields a statistical picture of the estimator's calibration: if the $(1-\alpha)$ credible intervals contain the true parameters in exactly a fraction $(1-\alpha)$ of trials, the estimator is well-calibrated.
\blue{Comparing the empirical coverage with the nominal level across a range of credibility thresholds produces a coverage-vs-credibility plot; departures from the diagonal indicate over- or underdispersion of the posterior, and the discrepancy is formally quantified through the empirical calibration error~\cite{Cranmer:2020,Hermans:2021rqv}.}

\blue{Coverage testing is} particularly well-suited to amortized SBI estimators, since evaluating the posterior for each new synthetic observation is generally fast (particularly for NPE) once the network has been trained.
\blue{For non-amortized methods, the computational cost of fitting a new posterior for each of the hundreds or thousands of required test datapoints would be prohibitive.}

\blue{Passing a coverage test is, however, a necessary but not a sufficient condition for proper calibration.
Because the test parameters are drawn from the prior, an estimator that ignores the data entirely and simply returns the prior passes the test perfectly; moreover, since the coverage is averaged over many observations, overconfident and underconfident regions of the parameter space can compensate each other~\cite{Hermans:2021rqv}.
Several complementary diagnostics are therefore often employed alongside it.
\emph{Simulation-based calibration} (SBC) sharpens the coverage test by requiring that the rank of each true parameter among the corresponding posterior samples be uniformly distributed, marginal by marginal~\cite{Talts:2018}.
\emph{Tests of accuracy with random points} (TARP) assess joint coverage through credible regions centered on random reference points, which extends the diagnostic to generative estimators whose density cannot be evaluated explicitly~\cite{Lemos:2023}.
Local diagnostics such as the \emph{local classifier two-sample test} (L-C2ST) train a classifier to distinguish true from estimated posterior samples for individual observations, exposing failures that the global averages above wash out~\cite{Linhart:2023}.
Most importantly, \emph{posterior predictive checks} validate the inference in data space: parameters drawn from the estimated posterior are fed back through the simulator, and the resulting synthetic observations are compared with the observed data.
A mismatch at this stage signals a badly wrong posterior or a misspecified simulator, a failure mode to which the parameter-space tests are blind. While possibly the most complete test, posterior predictive checks are also the most computationally expensive, since they require additional simulator calls for each posterior sample.}

In practice, many SBI analyses supplement these general calibration checks with problem-specific diagnostics.
% In Chapter~\ref{ch:6}, the internally estimated Bayesian uncertainties are cross-checked against frequentist error bars derived from the distribution of residuals over the validation set (see also Section~\ref{sec:3.1.4}).
% The close agreement between the two independent error estimates provides confidence in the reliability of the uncertainty quantification.
% \blue{In Chapter~\ref{ch:6}, the Bayesian SBI uncertainties are additionally cross-checked against an independent frequentist error estimate built from the validation-set residuals; the construction and outcome of this comparison are discussed in Section~\ref{sec:3.1.4}.}

